
% APPENDIX — ANALYTICAL GAUSSIAN MATRIX ELEMENTS


\chapter{Derivation of Gaussian Matrix Elements}
\label{app:gaussian-integrals}

For the hydrogenic Hamiltonian in atomic units,
\[
\hat H=-\frac12\nabla^2-\frac{Z}{r},
\]
we use normalised spherical Gaussian functions,
\[
\phi_i(r)=A_i e^{-\lambda_i r^2},
\qquad
p_{ij}=\lambda_i+\lambda_j.
\]
The matrix elements can be calculated analytically.
Since the functions are spherically symmetric, the
volume element becomes $d^3r=4\pi r^2\,dr$ after
integrating over the angles.

\section{Gaussian Normalisation}

Each Gaussian must satisfy
\[
\int|\phi_i|^2\,d^3r=1.
\]
Substituting the Gaussian function gives
\begin{align*}
1
&=4\pi A_i^2
  \int_0^\infty r^2e^{-2\lambda_i r^2}\,dr\\
&=4\pi A_i^2
  \frac{\sqrt{\pi}}{4(2\lambda_i)^{3/2}}\\
&=A_i^2
  \left(\frac{\pi}{2\lambda_i}\right)^{3/2},
\end{align*}
using
\[
\int_0^\infty r^2e^{-ar^2}\,dr
=\frac{\sqrt{\pi}}{4a^{3/2}}.
\]
Solving for $A_i$ gives
\begin{equation}
\boxed{
A_i=\left(\frac{2\lambda_i}{\pi}\right)^{3/4}.
}
\end{equation}

\section{Overlap Matrix}

The overlap between two Gaussian functions is
\begin{align*}
S_{ij}
&=\int\phi_i(\mathbf r)\phi_j(\mathbf r)\,d^3r\\
&=4\pi A_iA_j
  \int_0^\infty r^2e^{-p_{ij}r^2}\,dr\\
&=A_iA_j\frac{\pi^{3/2}}{p_{ij}^{3/2}}.
\end{align*}
Therefore,
\begin{equation}
\boxed{
S_{ij}=A_iA_j
\left(\frac{\pi}{p_{ij}}\right)^{3/2}.
}
\label{eq:appendix-overlap}
\end{equation}
The diagonal elements are $S_{ii}=1$ because each
Gaussian is normalised. The off-diagonal elements
are generally nonzero because the functions are
not orthogonal.

\section{Kinetic-Energy Matrix}

The kinetic-energy matrix elements are
\[
T_{ij}
=-\frac12\int\phi_i\nabla^2\phi_j\,d^3r.
\]
For a spherical Gaussian,
\begin{align*}
\frac{d\phi_j}{dr}
&=-2\lambda_jr\phi_j,\\
\nabla^2\phi_j
&=\frac1{r^2}\frac{d}{dr}
  \left(r^2\frac{d\phi_j}{dr}\right)\\
&=(4\lambda_j^2r^2-6\lambda_j)\phi_j.
\end{align*}
Substituting this into the kinetic-energy integral,
\begin{align*}
T_{ij}
&=3\lambda_jS_{ij}
 -2\lambda_j^2
  \int r^2\phi_i\phi_j\,d^3r.
\end{align*}
The remaining integral is
\begin{align*}
\int r^2\phi_i\phi_j\,d^3r
&=4\pi A_iA_j
  \int_0^\infty r^4e^{-p_{ij}r^2}\,dr\\
&=4\pi A_iA_j
  \frac{3\sqrt{\pi}}{8p_{ij}^{5/2}}\\
&=\frac{3}{2p_{ij}}S_{ij}.
\end{align*}
Substituting this result,
\begin{align*}
T_{ij}
&=3\lambda_jS_{ij}
  -\frac{3\lambda_j^2}{p_{ij}}S_{ij}\\
&=\frac{3\lambda_j(p_{ij}-\lambda_j)}
        {p_{ij}}S_{ij}\\
&=\frac{3\lambda_i\lambda_j}{p_{ij}}S_{ij}.
\end{align*}
Hence,
\begin{equation}
\boxed{
T_{ij}
=\frac{3\lambda_i\lambda_j}{p_{ij}}S_{ij}.
}
\label{eq:appendix-kinetic}
\end{equation}
This is symmetric in $i$ and $j$, as expected for
the kinetic-energy operator.

\section{Electron--Nucleus Potential Matrix}

The Coulomb potential is $V(r)=-Z/r$, giving
\begin{align*}
V_{ij}
&=-Z\int\frac{\phi_i\phi_j}{r}\,d^3r\\
&=-4\pi ZA_iA_j
  \int_0^\infty re^{-p_{ij}r^2}\,dr\\
&=-4\pi ZA_iA_j\frac{1}{2p_{ij}},
\end{align*}
using
\[
\int_0^\infty re^{-ar^2}\,dr=\frac1{2a}.
\]
Therefore,
\begin{equation}
\boxed{
V_{ij}=-\frac{2\pi ZA_iA_j}{p_{ij}}.
}
\label{eq:appendix-potential}
\end{equation}


\section{Variational Eigenvalue Problem}

The Hamiltonian matrix is obtained by adding the
kinetic and potential contributions:
\begin{equation}
    H_{ij}
    =\frac{3\lambda_i\lambda_j}{p_{ij}}S_{ij}
    -\frac{2\pi ZA_iA_j}{p_{ij}}.
    \label{eq:appendix-hamiltonian}
\end{equation}

Using the Gaussian expansion
$\psi_{\mathrm G}=\sum_{i=1}^{N}c_i\phi_i$,
the energy can be written as
\begin{equation*}
    E[\mathbf c]
    =\frac{\mathbf c^{\mathsf T}\mathbf H\mathbf c}
           {\mathbf c^{\mathsf T}\mathbf S\mathbf c}.
\end{equation*}

To find the coefficients that minimise the energy,
we use the normalisation condition
$\mathbf c^{\mathsf T}\mathbf S\mathbf c=1$
and introduce a Lagrange multiplier $E$:
\begin{equation*}
    \mathcal L(\mathbf c,E)
    =\mathbf c^{\mathsf T}\mathbf H\mathbf c
    -E(\mathbf c^{\mathsf T}\mathbf S\mathbf c-1).
\end{equation*}

Differentiating with respect to $\mathbf c$
and setting the result equal to zero,
\begin{align*}
    \frac{\partial\mathcal L}{\partial\mathbf c}
    &=2\mathbf H\mathbf c-2E\mathbf S\mathbf c=0,\\
    \mathbf H\mathbf c&=E\mathbf S\mathbf c.
\end{align*}

The lowest generalised eigenvalue gives the
approximate ground-state energy within the chosen
Gaussian basis. Since the matrix elements are
calculated analytically, numerical integration
is not needed to construct $\mathbf H$ or
$\mathbf S$.
